% SOURCE_STATEMENT: "Prove that, as the number of nuclei tends to infinity, the ground state of some neutral system of nuclei and electrons approaches a periodic limit, establishing the existence of crystals from quantum principles."
% Preparation and final self-audit: CONTRIBUTING.md prompts C and D; source/audit artifacts retained in the local physics_statements_2026_10_11 batch.
% !TeX program = lualatex
\documentclass[11pt]{article}
\usepackage[margin=25mm]{geometry}
\usepackage{fontspec}
\setmainfont{Latin Modern Roman}
\usepackage{amsmath,amssymb,mathtools}
\usepackage{unicode-math}
\setmathfont{Latin Modern Math}
\usepackage{xurl,hyperref}
\hypersetup{colorlinks=true,urlcolor=blue}
\setcounter{secnumdepth}{0}
\setlength{\parindent}{0pt}
\setlength{\parskip}{5pt}
\setlength{\emergencystretch}{3em}
\newcommand{\sref}[2]{\hyperlink{source:#1:#2}{[S#2]}}
\newcommand{\cataloguescope}[1]{\paragraph{Recorded target.}#1}
\newcommand{\R}{\mathbb R}
\newcommand{\C}{\mathbb C}
\newcommand{\N}{\mathbb N}
\newcommand{\Z}{\mathbb Z}
\newcommand{\Q}{\mathbb Q}
\newcommand{\F}{\mathbb F}
\newcommand{\eps}{\varepsilon}
\begin{document}
% UnsolvedMath ID 7900013; source catalogue code AMR-078-0013
\setcounter{section}{7900012}
\section{Existence of Quantum Crystals}\label{7900013}
{\small UnsolvedMath ID 7900013 · Mathematical Physics}
\subsection{Definitions and mathematical statement}

\paragraph{Shared notation.}$\mathbb N=\{1,2,\ldots\}$; the standard number fields have their usual meanings. All additional conventions are stated locally.


\paragraph{Definitions and notation.}A neutral quantum system has total nuclear charge equal to the number of electrons, in the corresponding charge units. A ground state minimizes the energy of its many-body Hamiltonian. A periodic thermodynamic limit would have particle-density or other specified observables invariant under a full-rank lattice of translations. A limit must specify normalization, boundary conditions and a convergence notion.

\paragraph{Catalogue wording retained for traceability.}
Prove that, as the number of nuclei tends to infinity, the ground state of some neutral system of nuclei and electrons approaches a periodic limit, establishing the existence of crystals from quantum principles.
\paragraph{Scope and source qualification.}The existential word some is preserved. The catalogue leaves the system, treatment of nuclei and topology of the periodic limit unspecified; a statement for every material or a prescribed lattice would be stronger.
\subsection{Short English statement}
Prove that a suitable neutral quantum system develops a periodic crystal structure as its size tends to infinity.
\subsection{Sources}
\begin{itemize}
\item[\hypertarget{source:7900013:1}{S1}] ULAM.AI and the UnsolvedMath Contributors, UnsolvedMath, record 7900013 (AMR-078-0013). Catalogue curation under CC BY 4.0; definitions and English exposition added locally.
\url{https://huggingface.co/datasets/ulamai/UnsolvedMath}.
\item[\hypertarget{source:7900013:2}{S2}] Barry Simon, Schrödinger operators in the twenty-first century, Mathematical Physics 2000, pp. 283–288, original Problem 13.
\url{https://web.ma.utexas.edu/mp_arc-bin/mpa?yn=00-78}.
\item[\hypertarget{source:7900013:3}{S3}] Wikipedia contributors, Simon problems, pinned revision 1335881350, The 2000 list. Existing catalogue wording adapted; CC BY-SA 4.0.
\url{https://en.wikipedia.org/w/index.php?title=Simon_problems&oldid=1335881350#The_2000_list}.
\end{itemize}
\label{end:7900013}
\end{document}
